% File: iel_lab02_cz.tex
% Source: https://github.com/josef-strnadel/howto-hints/blob/main/labs/IEL/iel_lab02_cz.tex
%
% Copyright (c) Josef Strnadel, https://www.fit.vut.cz/person/strnadel/

\documentclass[a4paper, 11pt]{report}
\usepackage[czech]{babel}
\usepackage{epsfig}
\usepackage{graphics}
\usepackage{multirow}
\usepackage{multicol}
\usepackage{fancyhdr}
\usepackage{lastpage}
\usepackage{latexsym}
\usepackage[hang, scriptsize, bf]{caption}

\usepackage[utf8]{inputenc}

\usepackage{enumerate}

\usepackage{tikz}
\usepackage[european]{circuitikz}

\usepackage{stackengine}
\usepackage{scalerel}
\usepackage{xcolor}
\newcommand\dangersign[1][2ex]{%
  \renewcommand\stacktype{L}%
  \scaleto{\stackon[1.3pt]{\color{red}$\triangle$}{\tiny\bfseries !}}{#1}%
}


%-------------------------------------------------------------------------------

\textwidth      170mm
\textheight     240mm
\hoffset        -10mm 
\voffset        -10mm
\oddsidemargin   5mm
\evensidemargin -5mm
\topmargin	    -5mm

%-------------------------------------------------------------------------------

\newcommand{\eduTitle}{Zadání laboratoře}
\newcommand{\eduID}{2}
\newcommand{\eduTopic}{Kondenzátor a RC článek}
\newcommand{\eduDetails}{
\footnotesize
Cíle: Experimentálně ověřit chování RC článku a porozumět významu RC článku pro praxi.
}
%
\newcommand{\subjIDlong}{Elektronika pro informační technologie}
\newcommand{\schoolDlong}{Vysoké učení technické v Brně}
\newcommand{\schoolDshort}{VUT v Brně}
\newcommand{\facultylDlong}{Fakulta informačních technologií}
\newcommand{\facultylDshort}{FIT}
%
\newcommand{\subjIDshort}{IEL}
%
\newcommand{\actYear}{2024} 
\newcommand{\acadYear}{2024/2025}

\usepackage[hidelinks]{hyperref}

%-------------------------------------------------------------------------------

\newcommand*\circled[1]{\tikz[baseline=(char.base)]{
            \node[shape=circle,draw,inner sep=2pt] (char) {#1};}}

%-------------------------------------------------------------------------------

\newcounter{cntInfo}
\newcommand{\info}[3]{\refstepcounter{cntInfo}
\paragraph*{
\circled{\thecntInfo}~{\sc \fbox{#1}} {\sc #2}  
} 
\paragraph{\textmd{#3}} }

%-------------------------------------------------------------------------------

\begin{document}

%-------------------------------------------------------------------------------

\pagestyle{fancy}
\renewcommand{\headrulewidth}{0pt}
\renewcommand{\theenumi}{\Alph{enumi}}  
\lhead{}
\chead{}
\rhead{}
\lfoot{\centering
\tiny \itshape \eduTitle ~č. \eduID \hspace{0.1mm} z předmětu \subjIDshort \hspace{1mm} 
(ak. r. \acadYear). \copyright~\actYear~Josef~Strnadel,~\facultylDshort~\schoolDshort. Připomínky zasílejte na \href{mailto:strnadel@fit.vut.cz}{strnadel@fit.vut.cz}\\ 
Časové razítko PDF dokumentu [\pdfcreationdate]. 
Sazba byla provedena systémem \LaTeX.
 }
\cfoot{}
%-------------------------------------------------------------------------------


%-------------------------------------------------------------------------------

\begin{center}
\scalebox{0.1}{\includegraphics{FIT_cernobile_CZ.pdf}}
\parbox{120mm}{
\textsc{\footnotesize\schoolDlong}, 
\textsc{\footnotesize\facultylDlong}\\
}
{
\Large
\textsc{\subjIDlong} (\textsc{\subjIDshort}), \textsc{ak. r. \acadYear}}

\hrulefill

\vspace{4mm}
\parbox{\linewidth}{
\centering
%
{\Huge \textsc{\eduTitle~č.~\fbox{\eduID}}}\\
{\textsc{,,\eduTopic''}}\\

}

\end{center}

{\it \eduDetails}

\hrulefill

%-------------------------------------------------------------------------------
\info{Motivace}{aneb ,,Proč tomu věnovat čas a jaké kompetence lze získat ?''}{
Na základě sady experimentů budete moci ověřit děje probíhající v RC článku při různých vstupních signálech,  
vliv dějů na odezvu RC článku a využití této odezvy v praxi.
}

%-------------------------------------------------------------------------------
\info{Výstup a způsob jeho hodnocení}{aneb ,,Co se ode mne očekává a co za to ?''}{
Za experimentální ověření dějů probíhajících v RC článku, 
jejich zaznamenání formou grafu 
a 
objasnění jejich využití v praxi lze získat až {\bf 3 body}.
}

%-------------------------------------------------------------------------------
\info{Prostředky}{aneb ,,Co je k dispozici ?''}{
~
\parbox{115mm}{
\vspace{-8mm}
Zdroj ss. napětí s omezením proudu,
nepájivé pole, 
krabička s prvky pro konstrukci obvodů (R, C, vodiče), 
generátor signálů (Obr. \thefigure; \dangersign[2ex]  \textcolor{red}{\textit{před zapnutím napájecího zdroje zkontrolujte, zda je generátor ke zdroji připojen správně, tj. vývod ``$+$'' k $5~V$, vývod ``$-$'' k $0~V$}} \dangersign[2ex]),
měřicí přístroje (multimetr, osciloskop). 
}
\hspace{0mm}
\parbox{45mm}{
\vspace{-8mm}
  \begin{center}
\scalebox{.8}{
\begin{tikzpicture}
\protect\draw
(0,0) -- (5,0) -- (5,-2.5) -- (3.5,-2.5) -- (3.5,-1.85) -- (1.5,-1.85) -- (1.5,-2.5) -- (0,-2.5) -- (0,0);
%
\protect\draw[line width = 2mm, color=red, opacity = .25] (0,0) -- (0.4,0) -- (0.4,-0.75) -- (0,-0.75) -- (0,.1);
\protect\node at(.2,-.4) {\bf $+$};
\protect\node at(1,-.4) {\textcolor{red!75}{\bf 5 V}};
%
\protect\draw[line width = 2mm, color=red, opacity = .25](4.3,0) -- (4.7,0) -- (4.7,-0.75) -- (4.3,-0.75) -- (4.3,.1);
\protect\node at(4.5,-.4) {\bf $+$};
\protect\node at(3.75,-.4) {\textcolor{red!75}{\bf 5 V}};
%
%
\protect\draw[line width = 2mm, color=blue, opacity = .25] (0.2,-2.25) -- (0.6,-2.25) -- (0.6,-1.5) -- (0.2,-1.5) -- (0.2,-2.35);
\protect\node at(0.4,-1.9) {\bf $-$};
\protect\node at(1.15,-1.55) {\textcolor{blue!75}{\bf 0 V}};
%
%
\protect\draw[line width = 2mm, color=blue, opacity = .25] (4.6,-2.25) -- (5.0,-2.25) -- (5.0,-1.5) -- (4.6,-1.5) -- (4.6,-2.35);
\protect\node at(4.8,-1.9) {\bf $-$};
\protect\node at(4.1,-2.2) {\textcolor{blue!75}{\bf 0 V}};
%
%
\protect\draw[very thick,fill=black,color=gray, opacity=.5] (3.65,-0.9) rectangle (4.25,-1.5);
\protect\filldraw[black] (3.95,-1.2) circle (5pt) node[left = 3mm]{\textcolor{black}{\bf SW1}};
%
%
%
\protect\draw[very thick,fill=blue,color=green, opacity=.25] (.75,-2.1) rectangle (1.5,-2.5);
\protect\node at(2.25,-2.3) {\scalebox{1}{\parbox{30mm}{\centering {\bf \textcolor{green!75!black}{výstup}}}}}; %\\ve %
\protect\draw(0.75,-2.1) -- (1.5,-2.1); 
\protect\draw(0.75,-2.1) -- (0.75,-2.5); 
\protect\node at(1,-2.35) {$\circ$};
\protect\node at(1.25,-2.35) {$\circ$};
%
\protect\draw[fill=none,color=gray!75, opacity=.75, dashed] (1.7,-0.2) rectangle (3.3,-1.75);
%
\end{tikzpicture}
  }
  \\{\bf \scriptsize Obrázek \thefigure:} {\scriptsize Generátor signálů}
  \end{center}
}
\vspace{-8mm}
}


\refstepcounter{figure}


%-------------------------------------------------------------------------------
\info{Základní schéma(ta) a vztahy}{aneb ,,Z čeho se bude vycházet ?''}{
~
}

\vspace{-6mm}
\begin{figure}[h!]
  \begin{center}
\scalebox{0.95}{
%
    \begin{circuitikz}
      \draw (1,2) node (in1) [ocirc]{} node[above] {\scriptsize 1}
     to[R,  l_=$R$, i_={}] (4,2)
      to[eC, l_=$C$, i_={}] (4,0) % The resistor
      to[short] (1,0)
      to[short] (1,0) node (in2) [ocirc]{}  node[above] {\scriptsize 2};

\draw node (A) at (4.5,0) {}
node (B) at (4.5,2) {}
(B) to[open, v^={}] (A);
\draw node (C) at (1.5,2.25) {}
node (D) at (3.5,2.25) {}
(C) to[open, v^={}] (D);
     \draw(0.5,0) node {\scriptsize a)};
     
\draw (4,2) to[short, *-o] (6,2) node[above] {\scriptsize 3};
\draw (4,0) to[short, *-o] (6,0) node[above] {\scriptsize 4};     

\node at(1,1) {$u_1(t)$};
\node at(6,1) {$u_2(t)$};
    \end{circuitikz}
} 
\hspace{8mm}
\scalebox{0.95}{
%
    \begin{circuitikz}
      \draw (1,2) node (in1) [ocirc]{} node[above] {\scriptsize 1}
     to[eC,  l_=$C$, i_={}] (4,2)
      to[R, l_=$R$, i_={}] (4,0) % The resistor
      to[short] (1,0)
      to[short] (1,0) node (in2) [ocirc]{}  node[above] {\scriptsize 2};

\draw node (A) at (4.5,0) {}
node (B) at (4.5,2) {}
(B) to[open, v^={}] (A);
\draw node (C) at (1.5,2.25) {}
node (D) at (3.5,2.25) {}
(C) to[open, v^={}] (D);
     \draw(0.5,0) node {\scriptsize b)};
     
\draw (4,2) to[short, *-o] (6,2) node[above] {\scriptsize 3};
\draw (4,0) to[short, *-o] (6,0) node[above] {\scriptsize 4};     

\node at(1,1) {$u_1(t)$};
\node at(6,1) {$u_2(t)$};
    \end{circuitikz}
} 

\vspace{-2mm}
   \caption{RC článek a) integrační, b) derivační.
   Základní veličiny a vztahy:
   $t$ je čas,
   $i(t)$ je proud procházející $R$, $C$,
   $u_1(t) = u_R(t) + u_C(t)$
   vstupní napěťový signál,  
   $u_2(t)$ je výstupní napěť. sig., přičemž  
   $u_2(t) = u_1(t) - u_R(t) = u_C(t)$ pro (a), 
$u_2(t) = u_1(t) - u_C(t) = u_R(t)$ pro (b).
Vztah mezi  $i(t)$, $u_C(t)$:
$u_c(t) = \frac{1}{C}\int_{}^{}i~dt$, 
$i(t) = C\cdot u'_C(t)$
pro $u_c(t)=0$ v $t=0$, elektrický náboj $Q$ na $C$, tj. $Q(t) = C \cdot u_C(t)$
a
proud $i(t)$ kondenzátorem $C$, tj. $i(t) = \frac{dQ(t)}{dt} = \frac{d(C\cdot u_C(t))}{dt} = C\cdot u'_C(t)$.
}
   \label{fig_rc}
  \end{center}
\end{figure}

\hspace{-6mm}
{

\vspace{-8mm}
Pojmy významné z hlediska i) dějů probíhajících v RC článku a ii)  využití RC článku v praxi: 

\begin{itemize}

\item
{\bf časová konstanta} ($\tau$ [$s$]) RC článku: 
\vspace{-5mm}
\begin{equation}
\tau = R\cdot C
\label{eq:tau}
\end{equation}

\vspace{-2mm}
\item
{\bf mezní úhlový kmitočet} ($\omega_0$ [$rad \cdot s^{-1}$]) RC článku: 
\vspace{-2mm}
\begin{equation}
\omega_{0} = \frac{1}{\tau} = 2\pi\cdot f_0
\label{eq:omega}
\end{equation}

\vspace{-5mm}
\item
{\bf napěťový přenos} ($\hat{A}_{u}(\omega)$ [$-$]) RC článku\footnote{$\hat{U}_1(j\omega)$, $\hat{U}_2(j\omega)$ jsou
frekvenčně (na $\omega$) závislé fázory příslušející $u_1(t)$, $u_2(t)$, resp. $A_{u_{dB}}(\omega)~[dB] = 20 \cdot \log_{10} A_u(\omega)$}: 
\vspace{-2mm}
\begin{equation}
\hat{A}_u(j\omega) = \frac{\hat{U}_2(j\omega)}{\hat{U}_1(j\omega)},~A_u(\omega) = |\hat{A}_u(j\omega)|
\label{eq:au}
\end{equation}


\end{itemize}



}


\newpage

%-------------------------------------------------------------------------------
\info{Postup samostatných činností}{aneb ,,Co dělat a na co si dát pozor ?''}{
~
}

\vspace{-7mm}
\hspace{-10mm}
\parbox{20mm}{
\flushleft
\dangersign[10ex] 
}
\parbox{148mm}{
\it
\textcolor{red}{
Používáme elektrolytické kondenzátory, u jejichž vývodů je nutno rozlišovat polaritu, aby nedošlo k poškození či dokonce k EXPLOZI kondenzátoru. Delší vývod je $+$, vývod~$-$ je na pouzdře kondenz. označen řetězcem symbolů "$-$" připomínajícím přerušovanou čáru.
}
}


\vspace{-4mm}

\begin{enumerate}[\bf {Experiment} 1:]
\item


\begin{enumerate}[i)]

\item 
Pro zvolené $R$, $C$ vypočtěte $\tau$ dle (\ref{eq:tau}) a {\bf zapojte} obvod dle Obr. \ref{fig_rc}a) či b).
Obvod {\bf rozšiřte} o přístroje pro
sledování
vývoje veličin $I$, $U_C$ (popř. $U_R$) v čase.

\item
Připojením ss./= napětí $u_1(t) = 5$ $V$ 
mezi svorky 1, 2 (viz Obr. \ref{fig_rc}) 
{\bf zahajte} nabíjení $C$
a 
{\bf sledujte}, 
i) jak se vyvíjejí veličiny $I$, $U_C$ (popř. $U_R$) v čase a 
ii)~jak změna $R$, $C$ (tj., $\tau$) ovlivňuje dynamiku dějů v obvodu 
({\it odpojením obvodu od zdroje napětí
můžete analogicky sledovat děje při vybíjení $C$}).


\end{enumerate}

\item

\begin{enumerate}[i)]
\item 
Pozorně {\bf sledujte} výklad vyučujícího k práci s osciloskopem\footnote{základní ovládání, informace o sondě, kalibračním vývodu, vizualizaci a měření v časové/hodnotové oblasti aj.}. 
Pomocí osciloskopu {\bf zjistěte} tvar signálu z kalibračního vývodu a 
{\bf odměřte} charakteristiky signálu v oblasti časové 
i 
napěťové. 

\item
{\bf Seznamte} se s obsluhou generátoru z Obr. 1 (dlouhý/krátký stisk SW1, výstup); 
pomocí osciloskopu {\bf analyzujte} jeho možnosti (tvar sig., frekv./napěť. rozsah).

\item 
{\bf Zapojte} jeden z obvodů z Obr. \ref{fig_rc}.
{\bf Nastavte a zapojte} generátor tak, 
aby na vstupu obvodu z Obr. \ref{fig_rc} produkoval napěťový signál $u_1(t)$ obdélníkového tvaru o zvolené frekvenci.
Pomocí osciloskopu {\bf sledujte} průběh $u_1(t)$  a 
(postupně také) průběh napětí na $R$, $C$. Průběhy napětí na $R$, $C$ {\bf objasněte
a nastiňte} jejich možné využití v praxi.

\end{enumerate}

\item
\begin{enumerate}[i)]
\item 
S ohledem na zjištěné vlastnosti generátoru z Obr. 1 pracujícího v ``sinusovém'' režimu
{\bf zvolte} $R$, $C$
tak, aby hodnota $f_0$ 
byla co nejblíže středu frekvenčního rozsahu generátoru pro tento režim. 

\item
{\bf Zapojte} jeden z obvodů z Obr. \ref{fig_rc}.
{\bf Nastavte} generátor tak, 
aby na vstupu obvodu generoval napěťový signál $u_1(t)$ sinusového tvaru o frekvenci $f = f_0$,
viz (\ref{eq:omega}),
{\bf změřte} 
amplitudy $u_1(t)$\footnote{tj. ${U}_1(\omega) = |\hat{U}_1(j\omega)|$},
$u_2(t)$\footnote{tj. ${U}_2(\omega) = |\hat{U}_2(j\omega)|$}
a {\bf spočítejte} modul (velikost) přenosu $A_u(\omega)$ dle (\ref{eq:au}). 
Analogicky {\bf postup opakujte} nejprve pro několik frekvencí $f~<~f_0$ až $f~\ll~f_0$,
poté 
pro několik frekvencí $f~>~f_0$ až $f \gg f_0$. 

\item
{\bf Vyneste} závislost $A_u(\omega)$ a $f$ do grafu
a ze závislosti {\bf vyvoďte}, zda se daný obvod chová jako filtr typu dolní či horní propust.

\end{enumerate}



\end{enumerate}

%-------------------------------------------------------------------------------
\info{Shrnutí, vyhodnocení a interpretace výsledků}{aneb ,,Jaká jsou zjištění ?''\vspace{-1mm}
}{
Experimentálně jste mohli ověřit, že kondenzátor se může nejen nabíjet, ale také (samo)vybíjet,
přičemž 
hodnota napětí na kondenzátoru je dána nejen 
 velikostí a směrem proudu protékajícího kondenzátorem, ale i výchozím nábojem (``pamětí'') v kondenzátoru.
Dynamika dějů v RC článcích je pak určena  hodnotou $\tau$,
čehož lze využít v řadě aplikací.
}

%-------------------------------------------------------------------------------
\vspace{-2mm}
\info{K zamyšlení/zapamatování}{aneb ,,Něco do dalšího studia a života.''\vspace{-1mm}
}{
Děje spojené s nabíjením/vybíjením kondenzátoru 
lze využít např. k odměřování času, detekci hran v obdélníkovém signálu, filtrování signálu, jeho derivaci či integraci.
Pro daný RC článek zkuste zjistit:
i) za jakých okolností bude velikost proudu obvodem maximální/minimální,
ii) za jakou dobu napětí na $C$ vzroste na určenou hodnotu při $u_1(t) = konst.$,   
iii) zda hodnota jeho $\tau$ umožňuje ``dobrou'' derivaci resp. integraci vstupního signálu,
iv) jaký je fázový rozdíl napětí a proudu na $C$. 
}


\end{document}
